\documentclass[11pt,a4paper]{article}
\usepackage{harmonikabau_preamble}

% == Dokumentmetadaten ====================
\newcommand{\hbdoknr}{0012}
\newcommand{\hbtitel}{Zungensteifigkeit}
\newcommand{\hbuntertitel}{Berechnung, Materialwahl und Einfluss auf Frequenz}
\newcommand{\hbheadtext}{Dok.\,0012 --- Zungensteifigkeit}
\newcommand{\hbfoottext}{Querverweise: Dok.\,0010}

\AtBeginDocument{%
  \fancyhead[L]{\small\hbheadtext}%
  \fancyhead[R]{\small\thepage}%
  \fancyfoot[C]{\small\color{hb@darkgray}\hbfoottext}%
}

\begin{document}
\hbmaketitle

{\small\itshape Die Steifigkeit der Stimmzunge bestimmt zwei Eigenschaften, die einander widersprechen: Ansprache und Stimmstabilität. Eine weiche Zunge spricht leichter an, verstimmt aber stärker mit dem Balgdruck. Eine steife Zunge hält den Ton stabiler, braucht aber mehr Druck zum Anschwingen.}

\section{Steifigkeit und Frequenz}

Frequenz und Steifigkeit hängen beide von der Geometrie ab, aber \textbf{unterschiedlich stark}:
\begin{equation}
f \propto \frac{t}{L^2}, \qquad k = \frac{E \cdot W \cdot t^3}{4 \cdot L^3}
\end{equation}
Die Steifigkeit $k$ hängt \textbf{kubisch} von $t$ ab, die Frequenz nur linear. 10\,\% dünner = 27\,\% weicher, aber nur 10\,\% tiefer.

\section{Verstimmung durch Balgdruck}

Der Bernoulli-Effekt erzeugt eine ``negative Feder'': $k_\text{eff} = k_\text{Zunge} - k_\text{Bernoulli}(p)$. Die Frequenzverschiebung: $\Delta f/f \approx -k_\text{Bernoulli}/(2 \cdot k_\text{Zunge})$.

\begin{keybox}
\textbf{Ergebnis:} Bass (k $\approx$ 80--150\,N/m): 4--11\,Cent/kPa (hörbar). Diskant (k $\approx$ 250--400\,N/m): 0,3--0,8\,Cent/kPa (kaum hörbar).
\end{keybox}

\section{Der Kompromiss}

$\Delta f \times p_\text{threshold} \approx \text{const}$ --- gute Ansprache und stabile Stimmung widersprechen sich physikalisch.

\section{Stimmgewicht vs. Verdünnung}

\textbf{Masse hinzufügen:} Senkt $f$, $k$ bleibt → gleiche Verstimmung. \textbf{Verdünnen:} Senkt $f$ und $k$ ($k \propto t^3$) → 37\,\% mehr Verstimmung bei 10\,\% Verdünnung.

\section{Zusammenfassung}
\begin{keybox}
1.~$k \propto t^3$: Kubische Kopplung. 2.~Verstimmung $\propto 1/k$. 3.~$\Delta f \times p_\text{th} \approx$ const. 4.~Stimmgewicht ist stimmstabiler als Verdünnung. 5.~Profilierung entkoppelt Masse und Steifigkeit.
\end{keybox}

\begin{center}
\textcolor{accentred}{\rule{0.6\textwidth}{1pt}}\\[4pt]
{\small\itshape Die Zunge sucht ihr Gleichgewicht zwischen Freiheit und Kontrolle --- genau wie der Spieler.}
\end{center}

\end{document}
